% =====================================================================
%  GridQ Platform Manual — shared preamble (FROZEN)
%  Every chapter is written against the macros defined here and must
%  compile with this preamble unchanged. Do not add packages in chapters.
% =====================================================================

% ---- Engine / encoding -------------------------------------------------
\usepackage[T1]{fontenc}
\usepackage[utf8]{inputenc}
\usepackage{lmodern}
\usepackage{microtype}
\usepackage[british]{babel}

% ---- Page geometry -----------------------------------------------------
\usepackage[a4paper,top=2.6cm,bottom=2.6cm,left=2.7cm,right=2.7cm,headheight=15pt]{geometry}
\usepackage{float}
\usepackage{lastpage}

% ---- Colour palette (brand) -------------------------------------------
\usepackage{xcolor}
\definecolor{gridqOrange}{HTML}{FF6B35}   % brand accent
\definecolor{gridqOrangeDk}{HTML}{C9521F}
\definecolor{gridqInk}{HTML}{0F172A}      % near-black text / headings
\definecolor{gridqSlate}{HTML}{334155}
\definecolor{gridqGrey}{HTML}{64748B}
\definecolor{gridqMist}{HTML}{E2E8F0}     % light rule / borders
\definecolor{gridqBg}{HTML}{F8FAFC}       % page-tint backgrounds
\definecolor{gridqBlue}{HTML}{2563EB}     % "rebuild" boxes / links
\definecolor{gridqBlueBg}{HTML}{EFF4FF}
\definecolor{gridqGreen}{HTML}{15803D}
\definecolor{gridqGreenBg}{HTML}{ECFDF3}
\definecolor{gridqAmber}{HTML}{B45309}
\definecolor{gridqAmberBg}{HTML}{FFFBEB}
\definecolor{gridqViolet}{HTML}{6D28D9}
\definecolor{gridqVioletBg}{HTML}{F5F3FF}

% Diagram node fills (kept muted so labels stay readable)
\definecolor{nodeProcess}{HTML}{E7EEF9}   % process / step
\definecolor{nodeProcessBd}{HTML}{2563EB}
\definecolor{nodeIO}{HTML}{FFF1E8}        % input / output data
\definecolor{nodeIOBd}{HTML}{FF6B35}
\definecolor{nodeDecision}{HTML}{FFF7E0}  % decision / gate
\definecolor{nodeDecisionBd}{HTML}{B45309}
\definecolor{nodeStore}{HTML}{ECFDF3}     % datastore / artifact
\definecolor{nodeStoreBd}{HTML}{15803D}
\definecolor{nodeTerm}{HTML}{EEF2F6}      % terminal start/end
\definecolor{nodeTermBd}{HTML}{334155}
\definecolor{nodeExt}{HTML}{F5F3FF}       % external / perimeter
\definecolor{nodeExtBd}{HTML}{6D28D9}

% ---- Maths / symbols (load before fonts that touch math) --------------
\usepackage{amsmath}
\usepackage{amssymb}
\usepackage{textcomp}

% ---- Tables ------------------------------------------------------------
\usepackage{booktabs}
\usepackage{array}
\usepackage{longtable}
\usepackage{tabularx}
\usepackage{makecell}
\usepackage{multirow}
\renewcommand{\arraystretch}{1.25}
% A left-aligned auto-width column for tabularx
\newcolumntype{L}{>{\raggedright\arraybackslash}X}
\newcolumntype{R}{>{\raggedleft\arraybackslash}X}

% ---- Lists -------------------------------------------------------------
\usepackage{enumitem}
\setlist{leftmargin=1.4em,itemsep=2pt,topsep=3pt}

% ---- Graphics ----------------------------------------------------------
\usepackage{graphicx}
\graphicspath{{figures/}{./}}

% ---- TikZ (flow diagrams) ---------------------------------------------
\usepackage{tikz}
\usetikzlibrary{positioning,shapes.geometric,arrows.meta,calc,fit,backgrounds,decorations.pathreplacing}

% Standard flowchart node styles — chapters use ONLY these.
\tikzset{
  fnbase/.style={
    font=\sffamily\small, align=center, text=gridqInk,
    minimum height=10mm, inner sep=3.2pt, line width=0.7pt,
    rounded corners=2pt
  },
  % process / pipeline step (rectangle)
  proc/.style={fnbase, rectangle, draw=nodeProcessBd, fill=nodeProcess,
    minimum width=34mm, text width=32mm},
  % wide process for top-level stages
  procwide/.style={proc, minimum width=46mm, text width=44mm},
  % input / output data (parallelogram)
  io/.style={fnbase, trapezium, trapezium left angle=70, trapezium right angle=110,
    draw=nodeIOBd, fill=nodeIO, minimum width=30mm, text width=30mm, trapezium stretches=true},
  % stored artifact / record (cylinder)
  store/.style={fnbase, cylinder, shape border rotate=90, aspect=0.18,
    draw=nodeStoreBd, fill=nodeStore, minimum width=30mm, text width=28mm,
    inner ysep=4pt},
  % decision / gate (diamond)
  decide/.style={fnbase, diamond, aspect=1.7, draw=nodeDecisionBd, fill=nodeDecision,
    minimum width=30mm, text width=26mm, inner sep=1pt},
  % external system / agentic perimeter (dashed)
  ext/.style={fnbase, rectangle, draw=nodeExtBd, fill=nodeExt, dashed,
    minimum width=32mm, text width=30mm},
  % terminal start / end (pill)
  term/.style={fnbase, rectangle, rounded corners=5mm, draw=nodeTermBd, fill=nodeTerm,
    minimum width=26mm, text width=24mm},
  % a faint label on the side
  note/.style={font=\sffamily\scriptsize\itshape, text=gridqGrey, align=left},
  % arrows
  flow/.style={-{Stealth[length=2.4mm,width=2mm]}, line width=0.9pt, draw=gridqSlate},
  flowback/.style={flow, dashed, draw=gridqGrey},
  flowlabel/.style={font=\sffamily\scriptsize, text=gridqSlate, fill=white, inner sep=1.5pt},
}
% Group-box behind a set of nodes (use on background layer)
\tikzset{
  groupbox/.style={rounded corners=4pt, draw=gridqMist, line width=1pt, fill=gridqBg},
}

% ---- Code / JSON listings ---------------------------------------------
\usepackage{listings}
\lstdefinestyle{gridqbase}{
  basicstyle=\ttfamily\footnotesize\color{gridqInk},
  backgroundcolor=\color{gridqBg},
  frame=single, rulecolor=\color{gridqMist}, framesep=6pt,
  xleftmargin=10pt, xrightmargin=4pt,
  breaklines=true, breakatwhitespace=false, columns=fullflexible,
  keepspaces=true, showstringspaces=false, upquote=true,
  tabsize=2, captionpos=b,
  commentstyle=\color{gridqGrey}\itshape,
  keywordstyle=\color{gridqBlue}\bfseries,
  stringstyle=\color{gridqGreen},
}
\lstdefinestyle{py}{style=gridqbase, language=Python,
  morekeywords={dataclass, frozen, Protocol, None, True, False}}
\lstdefinestyle{json}{style=gridqbase,
  morestring=[b]"}
\lstdefinestyle{txt}{style=gridqbase}
\lstset{style=gridqbase}
% inline code
\newcommand{\code}[1]{{\ttfamily\small\color{gridqInk}#1}}
\newcommand{\fpath}[1]{{\ttfamily\small\color{gridqSlate}#1}}

% ---- Callout boxes (tcolorbox) ----------------------------------------
\usepackage[most]{tcolorbox}
\tcbset{
  calloutbase/.style={
    enhanced, breakable, sharp corners=downhill, arc=2pt,
    boxrule=0pt, leftrule=3pt, top=6pt, bottom=6pt, left=9pt, right=9pt,
    fonttitle=\sffamily\bfseries\small, fontupper=\small,
    coltitle=gridqInk, before skip=10pt, after skip=10pt,
  }
}
% Key idea / takeaway
\newtcolorbox{keyidea}[1][Key idea]{calloutbase,
  colback=gridqAmberBg, colframe=gridqOrange, title={#1}}
% How to rebuild this (the recreate-from-scratch recipe)
\newtcolorbox{rebuild}[1][How to rebuild this step]{calloutbase,
  colback=gridqBlueBg, colframe=gridqBlue, title={#1}}
% Data contract: inputs / outputs
\newtcolorbox{contract}[1][Inputs and outputs]{calloutbase,
  colback=gridqGreenBg, colframe=gridqGreen, title={#1}}
% Plain-language definition / note
\newtcolorbox{plainnote}[1][In plain terms]{calloutbase,
  colback=gridqBg, colframe=gridqSlate, title={#1}}
% Caution / non-obvious rule
\newtcolorbox{caution}[1][Watch out]{calloutbase,
  colback=gridqVioletBg, colframe=gridqViolet, title={#1}}

% ---- Section / heading styling ----------------------------------------
\usepackage{titlesec}
\usepackage{titletoc}
\titleformat{\chapter}[display]
  {\sffamily\bfseries\color{gridqInk}}
  {\flushright\fontsize{52}{52}\selectfont\color{gridqMist}\thechapter}
  {-8pt}
  {\Huge\raggedright}
  [\vspace{2pt}{\color{gridqOrange}\titlerule[2.5pt]}]
\titlespacing*{\chapter}{0pt}{6pt}{22pt}
\titleformat{\section}
  {\sffamily\Large\bfseries\color{gridqInk}}{\thesection}{0.7em}{}
\titleformat{\subsection}
  {\sffamily\large\bfseries\color{gridqSlate}}{\thesubsection}{0.6em}{}
\titleformat{\subsubsection}
  {\sffamily\normalsize\bfseries\color{gridqSlate}}{}{0em}{}

% ---- Headers / footers -------------------------------------------------
\usepackage{fancyhdr}
\pagestyle{fancy}
\fancyhf{}
\renewcommand{\headrulewidth}{0.6pt}
\renewcommand{\footrulewidth}{0pt}
\renewcommand{\headrule}{\hbox to\headwidth{\color{gridqMist}\leaders\hrule height \headrulewidth\hfill}}
\fancyhead[L]{\sffamily\footnotesize\color{gridqGrey}\nouppercase{\leftmark}}
\fancyhead[R]{\sffamily\footnotesize\color{gridqGrey}GridQ Platform Manual}
\fancyfoot[C]{\sffamily\footnotesize\color{gridqGrey}\thepage}
\fancypagestyle{plain}{\fancyhf{}\renewcommand{\headrulewidth}{0pt}%
  \fancyfoot[C]{\sffamily\footnotesize\color{gridqGrey}\thepage}}

% ---- Captions ----------------------------------------------------------
\usepackage{caption}
\captionsetup{font={sf,small},labelfont={bf,color=gridqOrangeDk},
  labelsep=period, skip=8pt, width=0.92\textwidth}

% ---- Hyperlinks / PDF metadata ----------------------------------------
\usepackage[hidelinks]{hyperref}
\hypersetup{
  colorlinks=true, linkcolor=gridqOrangeDk, citecolor=gridqBlue,
  urlcolor=gridqBlue, filecolor=gridqBlue,
  bookmarksnumbered=true, bookmarksopen=true,
  pdftitle={GridQ Behind-the-Meter BESS Underwriting Platform — User Manual},
  pdfauthor={GridQ / Neura}, pdfsubject={Platform user and rebuild manual},
}
\usepackage{bookmark}

% ---- Small helpers shared by chapters ---------------------------------
% Module tag used at the head of a chapter, e.g. \modtag{M0}{Ingestion}{Perimeter}
\newcommand{\modtag}[3]{%
  \noindent\begin{tcolorbox}[enhanced,sharp corners,boxrule=0pt,colback=gridqInk,
    colupper=white,left=10pt,right=10pt,top=5pt,bottom=5pt,before skip=2pt,after skip=12pt]
  {\sffamily\bfseries\color{gridqOrange}#1}\;\;{\sffamily\bfseries\large #2}%
  \hfill{\sffamily\footnotesize\color{gridqMist}#3}\end{tcolorbox}}
% A figure float helper for flow diagrams (centres + frames lightly handled by caller)
% Convenience for "pipeline position" breadcrumb
\newcommand{\pipelinepos}[1]{\hfill{\sffamily\footnotesize\color{gridqGrey}Pipeline position: #1}\par\vspace{4pt}}

% Better \ttfamily line breaking for long identifiers/paths
\usepackage{seqsplit}

% Framed screenshot helper: \screenshot{<width fraction>}{<file>}
\newcommand{\screenshot}[2]{%
  {\setlength{\fboxsep}{0pt}\setlength{\fboxrule}{0.6pt}%
   \fcolorbox{gridqMist}{white}{\includegraphics[width=#1\textwidth]{#2}}}}
